\documentclass{article}
\usepackage{lipsum}

\usepackage{amssymb}            % For mathematical symbols
\usepackage{amsmath}            % For, say, math operators I guess
\usepackage[title]{appendix}    % For appendices. 
\usepackage{arydshln}          % For dashed lines

\usepackage{booktabs}           % For tables, mainly the \toprule \midrule \bottomrule
\usepackage{graphicx}           % Required for inserting images
\usepackage{hyperref}           % For hrefs
\usepackage{listings}           % For verbatim R code chunks
\usepackage{multirow}           % For, as you can guess, multirows
\usepackage{natbib}             % For bibliography
\usepackage{pdfpages}           % To include compiled pdf in project folder
\usepackage[caption = false]{subfig} % For floating figures
\usepackage{url}                % For bare url links
\usepackage[svgnames]{xcolor}   % For colouring text
\usepackage{xfrac}              % For inline fractions, advised by: https://tex.stackexchange.com/questions/128496

\newcounter{examplecounter}
\newenvironment{example}{\begin{quote}%
    \refstepcounter{examplecounter}%
  \textbf{Example \arabic{examplecounter}}%
  \quad
}{%
\end{quote}%
}
\begin{document}
\begin{example}\label{ex:simple}
  A simple example environment

  \lipsum[1]
\end{example}

Some text: \lipsum[4]

\begin{example}\label{ex:another}
  Another example
\end{example}


 %% -------- Unpublished  ---------------------------------- %%




\section{chunks}

$$
\ln \pi(z \vert \boldsymbol{\Sigma}) = C + \frac{1}{2}\sum_{i=1}^{n} d_i z_i ^{\top} \boldsymbol{\Sigma}^{-1} z_i -
\frac{1}{2}\sum_{i=1}^{n_z} \sum_{k=1}^{n_z}\rho w_{ik} z_i^{\top} \boldsymbol{\Sigma}^{-1} z_k
$$
%
Where $C$ is an additive constant independent on $z$. Let us introduce the row-wise vectorising operator $\mathrm{vec}_r$, which stacks row-wise the elements of a $n \times k$ matrix into a $kn \times 1$ vector. The first sum then reads as:
$$
 \sum_{i=1}^{n_z} d_i z^{\top}_i \boldsymbol{\Sigma}^{-1} z_i = 
\begin{pmatrix}
z_{11} \\ \vdots \\ z_{1k} \\ z_{21} \\ \vdots \\ z_{2k} \\ \vdots \\z_{n_z1} \\ \vdots \\ z_{n_zk}
\end{pmatrix}^{\top}
\begin{pmatrix}
d_1 \boldsymbol{\Sigma}^{-1} & 0 & \ldots & 0 \\
0 & d_2 \boldsymbol{\Sigma}^{-1} & \ldots & 0 \\
\vdots & \vdots & \ddots & \vdots \\
0 & 0 & \ldots & d_{n_z} \boldsymbol{\Sigma}^{-1}
\end{pmatrix}
\begin{pmatrix}
z_{11} \\ \vdots \\ z_{1k} \\ z_{21} \\ \vdots \\ z_{2k} \\ \vdots \\z_{n_z1} \\ \vdots \\ z_{n_zk}
\end{pmatrix}  = 
(\mathrm{vec}_rz) ^{\top} (\boldsymbol{D} \otimes \boldsymbol{\Sigma}^{-1}) \mathrm{vec}_r z
$$
And the second sum reads as:
$$
 \sum_{i=1}^{n_z} \sum_{k=1}^{n_z} \rho w_{ik} z^{\top}_i \boldsymbol{\Sigma}^{-1} z_k = 
\begin{pmatrix}
z_{11} \\ \vdots \\ z_{1k} \\ z_{21} \\ \vdots \\ z_{2k} \\ \vdots \\z_{n_z1} \\ \vdots \\ z_{n_zk}
\end{pmatrix}^{\top}\rho
\begin{pmatrix}
\boldsymbol{0_{k \times k}} & w_{12} \boldsymbol{\Sigma}^{-1} & \ldots & w_{1n_z} \boldsymbol{\Sigma}^{-1}\\
w_{21} \boldsymbol{\Sigma}^{-1} & \boldsymbol{0_{k \times k}} & \ldots & w_{2n_z} \boldsymbol{\Sigma}^{-1}\\
\vdots & \vdots & \ddots & \vdots \\
w_{n_z1} \boldsymbol{\Sigma}^{-1} & w_{n_z2} \boldsymbol{\Sigma}^{-1} & \ldots & \boldsymbol{0_{k \times k}}
\end{pmatrix}
\begin{pmatrix}
z_{11} \\ \vdots \\ z_{1k} \\ z_{21} \\ \vdots \\ z_{2k} \\ \vdots \\z_{n_z1} \\ \vdots \\ z_{n_zk}
\end{pmatrix}  = 
$$
$$
=(\mathrm{vec}_r z) ^{\top} \left( \rho \boldsymbol{W} \otimes \boldsymbol{\Sigma}^{-1} \right) \mathrm{vec}_r z
$$
\end{document}